\documentclass[11pt]{article}
\usepackage[T1]{fontenc}
\usepackage[utf8]{inputenc}
\usepackage{lmodern}
\usepackage{amsmath,amssymb,amsthm,mathtools}
\usepackage[margin=1in]{geometry}
\usepackage{microtype,booktabs,array,longtable}
\usepackage{xcolor}
\usepackage[colorlinks=true,linkcolor=blue!45!black,citecolor=blue!45!black,urlcolor=blue!45!black]{hyperref}
\usepackage{enumitem,fancyhdr}
\usepackage{listings}
\setlist{itemsep=3pt,topsep=5pt}
\setlength{\emergencystretch}{2em}
\setlength{\headheight}{14pt}
\pagestyle{fancy}
\fancyhf{}
\fancyhead[L]{\small A density-sensitive refinement of exact counting}
\fancyhead[R]{\small Research note}
\fancyfoot[C]{\thepage}
\newtheorem{theorem}{Theorem}[section]
\newtheorem{lemma}[theorem]{Lemma}
\newtheorem{proposition}[theorem]{Proposition}
\newtheorem{corollary}[theorem]{Corollary}
\theoremstyle{remark}
\newtheorem{remark}[theorem]{Remark}
\newcommand{\poly}{\operatorname{poly}}
\newcommand{\F}{\mathbb F}
\newcommand{\Z}{\mathbb Z}
\newcommand{\Q}{\mathbb Q}
\newcommand{\wt}{\operatorname{wt}}
\newcommand{\supp}{\operatorname{supp}}
\newcommand{\ceil}[1]{\left\lceil #1\right\rceil}
\newcommand{\floor}[1]{\left\lfloor #1\right\rfloor}
\newcommand{\upstream}{\href{https://github.com/openai/math/tree/42907c9e999cc00d6abe7a50bfc1023250edd1b9/preprints/The-Exponential-Time-Hypothesis-is-False-October-25-2026}{the pinned source snapshot}}
\setcounter{tocdepth}{1}
\makeatletter
\renewcommand*\l@section[2]{\@dottedtocline{1}{0em}{2em}{\bfseries #1}{\bfseries #2}}
\makeatother
\allowdisplaybreaks[2]
\hypersetup{pdftitle={A density-sensitive refinement of the October 2026 exact k-SAT counting algorithm},pdfauthor={ChatGPT, OpenAI},pdfsubject={Exact counting, expander codes, sparsity-sensitive tensor contraction}}
\title{\textbf{A density-sensitive refinement of the October 2026\\ exact \(\#k\)-SAT counting algorithm}}
\author{Research note prepared for Vladimir Reshetnikov\\
\small ChatGPT, OpenAI}
\date{26 October 2026}
\begin{document}
\maketitle
\begin{abstract}
We modify the parity-check reduction in OpenAI's preprint
\emph{The Exponential Time Hypothesis is False}, revised 25 October
2026. For fixed \(k\) and fixed \(A>0\), a \(k\)-CNF with \(n\)
variables and \(m\le Bn\) clauses, where
\(1\le B\le(\log n)^A\), can be counted deterministically in
\[
 (n+m)^{O_{k,A}(1)}
 2^{O_{k,A}\!\left(n\log(2+B)/\sqrt{\log n}\right)}
\]
bit operations and polynomial space, using the preprint's Boolean tensor
contraction theorem. The change replaces parity checks whose lengths
depend on \(\log m\) by checks whose lengths depend only polynomially on
\(B\log(2B)\). A refined accounting of the same incidence encoding then
uses \(O_k(n\log(2+B))\) Boolean wires.

For \(B=\Theta(\log\log n)\), the exponent improves as follows:
\[
 O_k\!\left(\frac{n\log\log n}{\sqrt{\log n}}\right)
 \quad\longrightarrow\quad
 O_k\!\left(\frac{n\log\log\log n}{\sqrt{\log n}}\right).
\]
This is a proved improvement for a growing-density family, not a reduction
of the unrestricted worst-case exponent. We also give exact improvements
to the contraction's Hermite precision budget, a sharpening of the
preprint's Max-\(k\)-SAT corollary, and explicit obstacles to several
tempting stronger claims. Reproducible finite checks accompany the note;
they are not a full implementation or a formalization of the new algorithm.
\end{abstract}
\begingroup\small\tableofcontents\endgroup
\newpage

\section{Result, scope, and comparison}\label{sec:result}
All logarithms in this note are base two unless written \(\ln\).
Their bases do not affect the displayed big-\(O\) classes. Set
\[
 \lambda=\ceil{\log n},\qquad B=\max\{1,\ceil{m/n}\}.
\]
Here \(n\) counts declared variables, including unused ones. Reading and
normalizing the input takes a polynomial factor in the original \(n+m\).
Normalization never increases the clause count.

The source is the version in commit
\texttt{42907c9e999cc00d6abe7a50bfc1023250edd1b9} of
\texttt{openai/math}. Theorems below use its Boolean tensor contraction
theorem and its exact parity-lift identities as established inputs
\cite{OAI26}. This note does not independently certify the entire
upstream proof or its Lean formalization.

\begin{theorem}[Density-sensitive exact counting]\label{thm:density}
Fix an integer \(k\ge2\) and a real constant \(A>0\).
For integers \(n\ge20\), \(B\ge1\), and \(m\le Bn\), with
\(B\le\lambda^A\), the number of satisfying assignments of a \(k\)-CNF
can be computed deterministically in
\begin{equation}\label{eq:density}
 (n+m)^{O_{k,A}(1)}
 2^{O_{k,A}\!\left(n\log(2+B)/\sqrt{\log n}\right)}
\end{equation}
bit operations and \((n+m)^{O_{k,A}(1)}\) bits of space.
The implied constants are uniform over the allowed values of \(B\).
\end{theorem}

A single algorithm also gives a useful all-input bound.

\begin{corollary}[A hybrid bound for every formula]\label{cor:hybrid}
For each fixed \(k\ge2\), there is a deterministic exact counter using
polynomial space whose time on every \(k\)-CNF is at most
\begin{equation}\label{eq:hybrid}
 (n+m)^{O_k(1)}
 2^{O_k\!\left(
 \frac{n}{\sqrt{\log n}}
 \min\{\log(2+B),\log(2+\lambda)\}\right)}.
\end{equation}
\end{corollary}
\begin{proof}
If \(B\le\lambda\), apply Theorem~\ref{thm:density} with \(A=1\).
Otherwise apply the source's arbitrary-density algorithm. At the cutoff,
\(\log(2+B)\le\log(2+\lambda)\). Beyond the cutoff the latter is
\(O(\log\log n)\), the numerator factor in the source's main theorem.
Bounded values of \(n\) are handled by exhaustive counting.
\end{proof}

\subsection{What is actually improved?}
The source already proves the arbitrary-density bound
\[
 (n+m)^{O_k(1)}
 2^{O_k(n\log\log n/\sqrt{\log n})}
\]
and a fixed-density bound with numerator \(n\).
Its older rational-interpolation argument has a
\((\log\log n)^2\) numerator; removing one of those two factors is
already done in its quantitative appendix and is not a contribution
of this note.

The source's direct clause-network proof is uniform in
\(s=n+m\). Combining it with the source's general algorithm yields,
for polylogarithmic \(B\), the exponent envelope
\begin{equation}\label{eq:old-envelope}
 O_k\!\left(\frac{n}{\sqrt{\log n}}
            \min\{1+B,\log\log n\}\right).
\end{equation}
The new density dependence is logarithmic in \(B\).
Table~\ref{tab:comparison} compares upper-bound scales, suppressing
polynomial factors and fixed constants.

\begin{table}[ht]
\centering
\small
\begin{tabular}{@{}p{.25\textwidth}p{.31\textwidth}p{.34\textwidth}@{}}
\toprule
Clause density \(B\) & Source exponent scale & Refined exponent scale\\
\midrule
\(O(1)\) &
\(n/\sqrt{\log n}\) &
\(n/\sqrt{\log n}\)\\[5pt]
\(\sqrt{\log\log n}\) &
\(n\sqrt{\log\log n}/\sqrt{\log n}\) &
\(n\log\log\log n/\sqrt{\log n}\)\\[5pt]
\(\Theta(\log\log n)\) &
\(n\log\log n/\sqrt{\log n}\) &
\(n\log\log\log n/\sqrt{\log n}\)\\[5pt]
\(2^{\sqrt{\log\log n}}\) &
\(n\log\log n/\sqrt{\log n}\) &
\(n\sqrt{\log\log n}/\sqrt{\log n}\)\\[5pt]
\((\log n)^\alpha\), fixed \(\alpha>0\) &
\(n\log\log n/\sqrt{\log n}\) &
Same order\\
\bottomrule
\end{tabular}
\caption{Comparison of proved exponent envelopes. These are not measured
running times or lower bounds on the source algorithm.}\label{tab:comparison}
\end{table}

More generally, the new scale is little-\(o\) of
\eqref{eq:old-envelope} whenever \(B\to\infty\) and
\(\log B=o(\log\log n)\). For arbitrary density the worst-case
numerator remains \(n\log\log n\). In particular, we do not claim
an \(O(n/\sqrt{\log n})\) exponent for every fixed-width CNF.

\section{Where the density improvement enters}\label{sec:mechanism}
For a formula with \(m\le Bn\) indexed clauses, let \(v_e(x)\in\F_2\)
be the violation bit of clause \(e\). A parity-check row \(a\in\F_2^m\)
gives the equation
\[
 \sum_{e=1}^m a_e v_e(x)=0.
\]
Each \(v_e\) has degree at most \(k\) and at most \(2^k\) monomials
over \(\F_2\). Every satisfying assignment passes every check.
The source first constructs checks that eliminate all sufficiently small
nonzero violation patterns, and then greedily adds further checks.

In its Lemma \texttt{sparse-checks}, the parameters are
\[
 \mu=\floor{\log m}+1,\quad
 P_{\rm src}=20000 B^2\mu^2,\quad
 D_{\rm src}=320000 B^2\mu^3.
\]
The initial rows have support at most \(D_{\rm src}\); any nonzero
accepted violation vector has weight at least \(m/P_{\rm src}\).
The construction uses an abelian Cayley graph obtained from small-bias
shifts and a deliberately coarse local-code distance estimate
\cite{OAI26}.

We replace those ingredients by a fixed-degree explicit expander,
raised to a power depending only on \(B\), and a binary-entropy estimate
for a high-rate local code. The resulting parameters satisfy
\begin{equation}\label{eq:new-parameters-summary}
 D=(B\log(2B))^{O(1)},\qquad
 P=O(B^2\log^2(2B)).
\end{equation}
They contain no \(\log m\) factor.

The later common-zero queries therefore have \(O(n)\) equations and
\[
 S=(B\log(2B))^{O_k(1)}
\]
monomials per equation. The source's incidence encoding, with its actual
sparsity retained, uses \(O_k(n\log(2+S))\) wires.
Since \(\log(2+S)=O_k(\log(2+B))\), its existing contraction theorem
gives \eqref{eq:density}. The next sections prove each part, including
construction time, rounding, and exactness.

\section{Initial parity checks depending only on density}\label{sec:checks}
\subsection{An explicit expander with controlled size}
We use the classical Margulis--Gabber--Galil family \cite{GG81}.
For every positive integer \(q\), its vertices are
\((\Z/q\Z)^2\). From \((x,y)\), the eight labeled transitions are
\[
 (x\pm2y,y),\quad (x\pm(2y+1),y),\quad
 (x,y\pm2x),\quad (x,y\pm(2x+1)),
\]
with arithmetic modulo \(q\).
They define a symmetric degree-eight multigraph. The nonconstant
eigenvalues of its normalized adjacency are at most
\(5\sqrt2/8<1\). An independently accessible formal proof of the
stronger absolute numerical-radius bound appears in
Karayel's \texttt{Expander\_Graphs\_MGG} development
\cite{Karayel23}, specifically \texttt{mgg\_numerical\_radius}.

Add eight labeled identity transitions. The resulting degree-sixteen
walk has normalized adjacency
\[
 H=\frac12(I+A/8),
\]
where \(A\) is the original adjacency. Its eigenvalues are nonnegative,
and its nonconstant eigenvalues are at most
\[
 \frac{8+5\sqrt2}{16}<\frac{19}{20}=:\rho.
\]
Thus an upper second-eigenvalue convention already suffices for our
construction. Every individual transition is a permutation.

Fix the integers \(B\ge1\) and
\begin{equation}\label{eq:code-parameters}
 w=\ceil{\log(2B)},\qquad
 \delta=\frac1{1024Bw},\qquad
 P=2(1024Bw)^2.
\end{equation}
Let \(t\) be the least nonnegative integer satisfying
\begin{equation}\label{eq:power-condition}
 (19/20)^t\le\delta/2,
 \qquad D=16^t.
\end{equation}
This condition is tested exactly as
\(2048Bw\,19^t\le20^t\). In particular \(t=O(\log(2B))\).
Since \((20/19)^{55}>16\),
\begin{equation}\label{eq:D-bound}
 32B\le D\le16(2048Bw)^{55}.
\end{equation}
For the lower bound, \(16^t\ge(20/19)^t\ge2048Bw\).
The upper bound follows from minimality of \(t\).

If \(m>D\), choose
\[
 q=\ceil{\sqrt{m/D}},\quad Q=q^2,\quad d=D.
\]
Take the \(t\)-step lazy walk on \(Q\) vertices and its bipartite
double cover. There are \(Q\) vertices on each side, degree \(d\),
and \(E=Qd\) distinct labeled edges. Walks with the same endpoints
are still distinct edges. The rounding gives
\begin{equation}\label{eq:rounding}
 m\le E<4m.
\end{equation}
The normalized second singular value of the biadjacency matrix is at
most \(\epsilon=\rho^t\le\delta/2\).

If \(m\le D\), use one vertex on each side with \(d=m\) parallel edges,
so \(Q=1\), \(E=m\), and \(\epsilon=0\).
This avoids padding a small input to \(D\) edges.
At either endpoint the incident ports receive any fixed ordering.
A length-\(t\) word in the sixteen transitions identifies a permutation,
so the endpoint and its port can be computed using polynomially many
bit operations in \(t+\log q\).
For both cases, the mixing inequality is
\begin{equation}\label{eq:mixing}
 e(U,V)\le d\frac{|U||V|}{Q}
              +\epsilon d\sqrt{|U||V|}.
\end{equation}
It follows by separating the constant and orthogonal parts of the two
indicator vectors and applying the singular-value bound.
For \(Q=1\) it holds directly with \(\epsilon=0\).

\subsection{A local code with entropy-controlled distance}
\begin{lemma}\label{lem:local-code}
For every \(d\ge1\) there is a binary matrix with
\(c=\ceil{d/(32B)}\) rows and \(d\) columns such that every nonzero
kernel word has weight at least \(\delta d\).
It can be found deterministically in \(2^{O(d^2)}\) time and
polynomial space in \(d+B\).
\end{lemma}
\begin{proof}
Write \(H_2(z)=-z\log z-(1-z)\log(1-z)\).
For \(0<z<1\), the elementary entropy estimate gives
\(H_2(z)\le z\log(\mathrm e/z)\).
Since \(\log B\le w-1\), \(\log w\le w-1\), and
\(\log(1024\mathrm e)<12\), we have
\[
 H_2(\delta)
 \le\frac{\log(1024\mathrm e Bw)}{1024Bw}
 \le\frac{12}{1024B}<\frac1{64B}.
\]
The number of binary words of positive weight below \(\delta d\)
is at most \(2^{dH_2(\delta)}\). For a uniformly random
\(c\times d\) binary matrix, each fixed nonzero word is in the
kernel with probability \(2^{-c}\). A union bound gives failure
probability less than
\[
 2^{d/(64B)-c}\le2^{-d/(64B)}<1.
\]
If \(\delta d\le1\), there are no forbidden nonzero words and the
claim is immediate.

Enumerate matrices lexicographically and test all words of weight
strictly less than \(\delta d\), retaining the first matrix with no
such nonzero kernel word. There are at most \(2^{cd}\) matrices and
\(2^d\) test words per matrix. The current matrix, word, and syndrome
require only polynomial space. The probabilistic argument proves
that this deterministic finite search terminates.
\end{proof}

\subsection{Global distance, shortening, and the row budget}
At every vertex of the bipartite graph impose the local code's \(c\)
checks on its incident edge bits. Keep any \(m\) edges and fix all other
edge coordinates to zero. Restricting the local checks to those \(m\)
coordinates produces a binary matrix \(A^{(0)}\).

\begin{lemma}[Density-sensitive initial checks]\label{lem:checks}
For \(1\le m\le Bn\), the matrix \(A^{(0)}\) has at most \(3n\)
rows, each supported on at most \(D\) coordinates, and
\[
 0\ne z\in\ker_{\F_2} A^{(0)}
 \quad\Longrightarrow\quad
 \wt(z)\ge m/P.
\]
Its construction costs a polynomial in \(n+m\) times
\(2^{O((Bw)^{110})}\), and uses polynomial space in \(n+m+B\).
\end{lemma}
\begin{proof}
First consider a nonzero accepted word on all \(E\) edges.
Let its weight be \(h\), with touched vertex sets \(U,V\).
The local distance implies
\[
 h\ge\delta d\max\{|U|,|V|\}
    \ge\delta d\sqrt{|U||V|}.
\]
All selected edges run between \(U\) and \(V\). Combining this with
\eqref{eq:mixing} and dividing by \(d\sqrt{|U||V|}>0\) yields
\[
 \sqrt{|U||V|}\ge(\delta-\epsilon)Q.
\]
Consequently,
\[
 h\ge\delta(\delta-\epsilon)Qd
       \ge \frac{\delta^2 E}{2}
       \ge \frac{m}{P}.
\]
The argument also applies when \(Q=1\), with \(\epsilon=0\).
A word accepted after shortening extends by zeros to an accepted word
on all \(E\) edges, so shortening preserves this distance bound.

If \(m>D\), use \(d=D\ge32B\) and \eqref{eq:rounding}:
\[
 2Qc\le\frac{Qd}{16B}+2Q
       \le\frac{Qd}{8B}
       <\frac{m}{2B}\le\frac n2.
\]
If \(m\le D\), then \(Q=1,d=m\), and
\[
 2c\le\frac{m}{16B}+2\le\frac n{16}+2\le3n.
\]
Every row reads at most \(d\le D\) edges.

The graph and its \(O(m)\) labeled edges can be listed in polynomial
time in their description parameters. The local-code search costs
\(2^{O(d^2)}\le2^{O(D^2)}\).
Equation~\eqref{eq:D-bound} gives the stated explicit exponent.
The full check matrix and all streamed searches use polynomial space.
\end{proof}

\begin{remark}[Asymptotic rather than practical construction]
The exponent \(110\) is deliberately crude, coming from a simple
explicit expander and exhaustive local-code construction.
For fixed polylogarithmic \(B\), it contributes only
\(2^{(\log n)^{O(1)}}\), asymptotically below the main running time.
It can nevertheless make the finite threshold extremely large.
No practical speedup is inferred from this construction.
\end{remark}

\section{A sparse-query bound retaining the actual sparsity}\label{sec:query}
The next lemma refines the accounting of the source's
\texttt{tensor-sparse-query} lemma. It introduces no new interpolation
or contraction identity.

\begin{lemma}[Sparsity-sensitive common-zero counting]\label{lem:query}
Fix positive constants \(c_0,c_1\) and an integer \(k\ge1\).
Suppose \(a\le c_0n\) polynomials in \(n\) Boolean variables over
\(\F_2\) each have degree at most \(k\) and at most \(S\) listed
monomials, where \(1\le S\le\lambda^{c_1}\).
Their exact number of common zeros is computable in
\[
 n^{O_{k,c_0,c_1}(1)}
 2^{O_{k,c_0,c_1}\!\left(
 n\log(2+S)/\sqrt{\log n}\right)}
\]
time and polynomial space.
\end{lemma}
\begin{proof}
Normalize powers and cancel duplicate monomials modulo two.
Let \(N_0\) be the number of test--variable incidences, counting a
variable once in each test's scope. Then \(N_0\le akS\).
Fix, in all possible ways, variables belonging to more than
\[
 D_*=\max\{1,\ceil{N_0/n}\}\ceil{\sqrt\lambda}
\]
scopes. Fewer than \(n/\sqrt\lambda\) variables are fixed, so
there are at most \(2^{n/\sqrt\lambda}\) branches.
Each remaining variable appears in at most
\(O_{k,c_0}(S\sqrt\lambda)\) scopes.
Constant equations and empty systems are handled immediately.

On a branch keep \(n\) as the original variable-count parameter.
Use the source's parity lifts over
\[
 \mathcal B=(\Z/2^{n+1}\Z)[V]/(V^{n+1}).
\]
Its total-order identity implies that, after multilinearizing in the
separate incidence variables \(z_{iv}\), the retained product \(F(z,V)\)
has total incidence degree at most
\begin{equation}\label{eq:incidence-K}
 K=k(a+2n)=O_{k,c_0}(n).
\end{equation}
If \(\mathcal S\) sums the \(V\)-coefficients after multiplication in
the truncated ring, the branch count is congruent modulo \(2^{n+1}\) to
\[
 \mathcal S\!\left(\sum_{x\in\{0,1\}^{\text{free variables}}}
              F((z_{iv}=x_v),V)\right).
\]
The truncation precedes coefficient summation; treating
\(\mathcal S\) as a ring homomorphism would be incorrect.

Partition the tests and free variables separately into
\(J=\ceil{n^{1/4}}\) groups. On either side a group touches at most
\[
 O_{k,c_0}(n^{3/4}S\sqrt\lambda)
\]
incidences. For each group pair \(g,h\), let \(N_{gh}\) be its incidence
count. Enumerate degree arrays \(d_{gh}\), with
\(0\le d_{gh}\le N_{gh}\) and \(\sum d_{gh}\le K\).
There are at most
\[
 (K+1)^{J^2}=2^{O_{k,c_0}(\sqrt n\log n)}
\]
arrays. For one array, encode the chosen \(d_{gh}\)-subset of the
\(N_{gh}\) incidences by its lexicographic rank, using
\(\ceil{\log\binom{N_{gh}}{d_{gh}}}\) Boolean wires.
Invalid rank strings receive tensor entry zero.
Writing \(N=\sum N_{gh}\le N_0\), the wire count obeys
\begin{align}
 M&\le J^2+\log\prod_{g,h}\binom{N_{gh}}{d_{gh}}\notag\\
  &\le J^2+\log\binom{N+K}{K}\notag\\
  &\le J^2+K\log\!\left(\frac{\mathrm e(N+K)}K\right)
     =O_{k,c_0}(n\log(2+S)).\label{eq:wire-budget}
\end{align}
This last step is the quantitative change from the source:
we retain \(S\), instead of replacing it immediately by
\((\log n)^{O(1)}\).

For completeness, the test-group tensor entry is the product of the
coefficients of its specified incidence subsets in the lifted factors.
If \(S_i\) is a selected subset of test \(i\)'s scope, that coefficient
can be evaluated by the finite-difference identity
\[
 [z_{i,S_i}]\overline B_i(z_i,V)
  =\sum_{T\subseteq S_i}(-1)^{|S_i|-|T|}
                     B_i(\mathbf1_T,V).
\]
Its cost is \(2^{|S_i|}\le2^{kS}\) times a polynomial in \(n\).
The variable-group tensor entry is \(2^{u_h}\) modulo \(2^{n+1}\),
where \(u_h\) is the number of free variables in that group untouched
by the selected incidences. An untouched variable contributes two
choices; a touched one must be one in the monomial sum.
This includes variables absent from every test.
These entries and the source's exact identities show that the sum
of the network contractions is the required branch residue.

Apply the source's Boolean tensor contraction theorem with
\[
 s=\ceil{n\log(2+S)}.
\]
Every hypothesis holds uniformly for sufficiently large \(n\):
\begin{itemize}
\item Equation~\eqref{eq:wire-budget} gives \(O_{k,c_0}(s)\) wires.
There are \(2J\) vertices and polynomial-size descriptions.
\item The maximum degree is
\(O_{k,c_0}(n^{3/4}S\sqrt\lambda)\), at most \(s^{9/10}\)
because \(S\le\lambda^{c_1}\).
\item Both the root precision and the truncation order are \(n+1=O(s)\).
\item An entry can be evaluated in polynomial space and time
\(n^{O(1)}2^{kS}\le n^{O(1)}2^{O(s^{9/10})}\).
\end{itemize}
Also \(\log s=\Theta_{c_1}(\log n)\). Each contraction thus costs
\[
 n^{O(1)}2^{O(n\log(2+S)/\sqrt{\log n})}.
\]
The variable-fixing and degree-array exponents are absorbed in this
bound. Process branches and arrays sequentially and add their counts
modulo \(2^{n+1}\). The final least nonnegative residue is the integer
answer, since it lies between zero and \(2^n\).
The finite exceptional range can be handled by direct assignment
enumeration; its threshold depends only on the fixed constants.
\end{proof}

\section{Completing the exact-counting algorithm}\label{sec:algorithm}
\subsection{Greedy completion of the checks}
Let \(N(A)\) be the number of assignments passing all parity checks
in the matrix \(A\). Write
\[
 N(A)=\#F+R(A),
\]
where \(R(A)\) counts the nonsatisfying assignments still passing.
The algorithm does not know \(R(A)\) or \(\#F\) in advance.

Starting with \(A^{(0)}\), enumerate candidate rows by all ordered
lists
\[
 ((e_1,c_1),\ldots,(e_P,c_P))
       \in([m]\times\{0,1\})^P.
\]
Repeated indices are allowed; their coefficient bits add modulo two.
By Lemma~\ref{lem:checks}, a nonsatisfying assignment passing
\(A^{(0)}\) has at least \(m/P\) violated clauses.
A uniform list misses this support with probability at most
\((1-1/P)^P\le e^{-1}<1/2\). Conditional on a hit, one independently
uniform coefficient bit makes the parity uniform.
Thus at least one quarter of candidate lists reject each such
assignment.

At each stage use Lemma~\ref{lem:query} to evaluate \(N(A;h)\)
for every candidate row \(h\), and append a minimizer.
The unknown summand \(\#F\) is the same for every candidate,
so this minimizes \(R(A;h)\) as well. Averaging gives
\[
 R(A;h_{\rm chosen})\le\tfrac34 R(A).
\]
After \(3n\) rounds,
\[
 R(A)\le2^n(3/4)^{3n}=(27/32)^n<1.
\]
It is an integer, hence zero. A final common-zero count returns \(\#F\).
This is the source's deterministic greedy argument with the improved
initial-check parameters; no random choice remains.

\subsection{Time, space, and uniformity}
Every query has at most \(6n\) equations. Its monomial sparsity is bounded by
\[
 S=2^k\max\{D,P\}=(B\log(2B))^{O_k(1)}.
\]
For fixed \(A\) and \(B\le\lambda^A\), this is at most a fixed power
of \(\lambda\), uniformly in \(B\), for sufficiently large \(n\).
Moreover,
\[
 \log(2+S)=O_k(\log(2+B)).
\]
There are at most \(3n(2m)^P+1\) calls to the common-zero counter.
Before absorbing lower-order factors, the exponential work is at most
\(2^{O_{k,A}(E(n,B,m))}\), where
\begin{align}
 E(n,B,m)={}&
 \frac{n\log(2+B)}{\sqrt{\log n}}+(B\log(2B))^{110}\notag\\
 &+B^2\log^2(2B)\log(2m).\label{eq:full-overhead}
\end{align}
Since \(m\le n\lambda^A\), both added terms are fixed powers of
\(\log n\). They are \(o(n/\sqrt{\log n})\), uniformly over the
promised range. This proves the time assertion of
Theorem~\ref{thm:density}. The explicit overhead is included to make clear
why exponentially many local-code candidates do not invalidate the
asymptotic result.

All searches and calls are sequential. The graph, current local code,
tested word, parity rows, candidate list, and exact counts have
polynomial-size descriptions. Each query has polynomial space by
Lemma~\ref{lem:query}. The current output count uses only \(n+1\)
bits. This proves the space assertion.

\subsection{Algorithm specification}
The following specifies the proposed hybrid algorithm.
\begin{enumerate}
\item Normalize clauses. Return zero for an empty clause, and \(2^n\)
for an empty formula. Compute \(B=\max\{1,\ceil{m/n}\}\) using the
remaining clauses.
\item If \(B>\lambda\), invoke the source's arbitrary-density counter.
Otherwise compute \(w,\delta,P,t,D\) from
\eqref{eq:code-parameters}--\eqref{eq:power-condition}.
\item Construct the graph, local code, and shortened initial matrix
\(A^{(0)}\) from Section~\ref{sec:checks}.
\item For \(3n\) rounds, enumerate all candidate lists, count each
augmented system using the sparsity-sensitive routine, and append the
lexicographically first minimizing candidate.
\item Count the final system exactly and return its count.
\end{enumerate}
This is a complete algorithmic description given the source contraction
routine. The accompanying programs validate component identities and
small instances; they do not implement the full contraction routine.

\section{A separate improvement of the contraction constants}\label{sec:retuning}
The density theorem above is independent of the following retuning.
These changes also apply to the source's general tensor contraction,
but do not change the order \(s/\sqrt{\log s}\).
In this section only, \(\lambda=\ceil{\log s}\), where \(s\) is the
source contraction size parameter.

\subsection{Use all Hermite conditions}
At an active call, retain the source's grid
\(S_0=\{-L,\ldots,L-1\}\), where \(L=2^a\).
For \(i=v_2(t)<u\), set \(q_i=h(u-i)\); all higher levels get
multiplicity zero. There are \(L/2^i\) points of valuation \(i\).
Consequently the exact number of Hermite conditions is
\begin{equation}\label{eq:exact-conditions}
 h\sum_{i=0}^{u-1}(L/2^i)(u-i)
   =hL(2u-2+2^{1-u}).
\end{equation}
The source uses only the odd points to certify enough conditions and
chooses \(h=\ceil{3M/u}\).
For \(u\ge2\), \(h=\ceil{M/(u-1)}\) already makes
\eqref{eq:exact-conditions} strictly greater than
\(M(2L-1)\), the degree bound.
The exact smallest integral slope is
\begin{equation}\label{eq:exact-h}
 h=\floor{\frac{M(2L-1)}
     {\sum_{i=0}^{u-1}(L/2^i)(u-i)}}+1.
\end{equation}
It also handles \(u=1\).

\subsection{Use the exact valuation loss}
Let \(D_i=L/2^i\), \(\alpha=2^{a+1}\), and let \(p\) be the
call's input precision. The source's Hermite weight \(E_t\) satisfies
\[
 v_2(E_t(0))
 =D_i\sum_{j=1}^{a-i}\frac j{2^j}(q_{i+j}-q_i).
\]
For the linear taper, its exact negative valuation is
\begin{equation}\label{eq:loss}
 \Lambda_i
  =h\sum_{j=1}^{a-i}j\min\{j,u-i\}\,2^{a-i-j}
  \le6hD_i.
\end{equation}
The inverse of \(E_t(t+\alpha z)\) has integral coefficients.
Evaluation of a jet of degree below \(q_i\) at \(-t/\alpha\)
loses at most \((q_i-1)(a+1-i)\) further powers of two.
Therefore it suffices to take
\begin{equation}\label{eq:exact-precision}
 p_t=p+\Lambda_i+(q_i-1)(a+1-i).
\end{equation}
Indeed, an error of valuation at least \(p_t\) in any jet coefficient
still has valuation at least \(p\) after these two losses.
Adding the errors does not reduce their minimum valuation.

The simpler replacement
\(p_t=p+6hD_i+(q_i-1)(a+1-i)\) is also valid.
The source uses \(p+8hD_i+8M(a+1)\).
Its inherited-precision argument still applies, and uniformly over
active calls the simplified new budget is
\[
 p_t=(6+o(1))MD_i/u,
\]
in place of \((24+o(1))MD_i/u\) for the source's stated choices.

\subsection{Retain more valuation levels}
Keep \(r=64\) groups, maximum degree \(\Delta=\ceil{s^{9/10}}\),
and the source leaf cutoff \(H=\ceil{s/\floor{\sqrt\lambda}}\).
Use
\[
 u_d=\floor{\eta\lambda/2^d},\qquad
 a_d=2\sum_{j=0}^d u_j.
\]
The source sets \(\eta=1/131072\).
The required collision-packing condition permits every fixed
\(\eta<1/960\), because
\[
 L=O(s^{4\eta}),\quad D_{\rm arr}=O(s^{48\eta}),\quad
 D_{\rm arr}^2\Delta=O(s^{9/10+96\eta})=o(H).
\]
The other two packing inequalities also hold:
\(LD_{\rm arr}=O(s^{52\eta})=o(H)\) and \(2L\le\Delta\)
eventually. A conservative replacement is
\(\eta=1/2048\), for which
\[
 4\eta=\frac1{512},\quad48\eta=\frac3{128},\quad
 \frac9{10}+96\eta=\frac{303}{320}<1.
\]
Tensor-entry generation still costs
\(\exp(O(s^{1/512}\log s))=\exp(o(s^{9/10}))\).
Interpolation uses fewer than \(s^4\) points eventually, and
all precisions and descriptions remain polynomial in \(s\).
The source's convenient bound \(L,D_{\rm arr}\le s^{1/1000}\)
must be replaced by these valid bounds; it is not retained unchanged.

\subsection{Meaning of the constant improvement}
There are at most \(6M/D_i\) modes of each positive weight.
Optimizing the same generating-function bound used in the source gives
\[
 \ln N_{\rm modes}
 \le\inf_{\theta>0}
       \left(\theta p_t+\frac{\pi^2M}{D_i\theta}\right)
 =2\pi\sqrt{Mp_t/D_i}.
\]
The precision change halves its leading coefficient.
Increasing \(\eta\) by \(64\) contributes a further factor \(8\).
Thus the matched \emph{proved mode-enumeration coefficient} decreases
by a factor \(16\).

For clarity, using the exact child shrink bound \(3/32\), an explicit
base-two coefficient per root edge constant \(C\) becomes
\[
 A_{\rm new}
 =\frac{2\pi\sqrt6}
        {\ln2\,\sqrt{1/2048}\,(1-3\sqrt2/32)}
 \simeq1158.422114.
\]
The same optimized analysis of the source's choices gives
\(A_{\rm old}\simeq18534.753832=16A_{\rm new}\).
The separate leaf contribution has coefficient one and is unchanged:
\[
 T\le\poly(s+\mathcal L)\,
 2^{(1+C A_{\rm new}+o(1))s/\sqrt{\log s}}.
\]
This compares upper bounds, not measured runtimes or a proved ratio
between the actual runtimes. It does not improve the power \(1/2\)
of the logarithm.

\section{A sharper \texorpdfstring{Max-\(k\)-SAT}{Max-k-SAT} consequence}\label{sec:maxsat}
The source's additive-\(\varepsilon\) approximation corollary states an
exponent \(O_{k,\varepsilon}(n(\log\log n)^2/\sqrt{\log n})\).
Its proof samples \(O(n/\varepsilon^2)\) clauses, maximizes their
satisfied fraction through a linear-size threshold circuit, and uses
a counting-preserving circuit-to-\(3\)-CNF encoding
\cite{OAI26}. For fixed \(k,\varepsilon\), both the number of variables
and the number of clauses in that encoding are \(O_{k,\varepsilon}(n)\).
Apply the fixed-density part of the source's main theorem directly.

\begin{corollary}\label{cor:maxsat}
For fixed \(k\ge2\) and \(\varepsilon>0\), an assignment satisfying at
least \(\operatorname{OPT}(F)-\varepsilon m\) clauses can be found with
probability \(1-2^{-\Omega(n)}\) in
\[
 (n+m)^{O_{k,\varepsilon}(1)}
 2^{O_{k,\varepsilon}(n/\sqrt{\log n})}
\]
time and polynomial space.
\end{corollary}
\begin{proof}
The source's uniform sampling estimate, threshold circuit, binary
search, and self-reduction are unchanged.
Each resulting \(3\)-CNF has fixed clause density, so the
fixed-density counting bound applies. Polynomially many calls and
the source's bounded sampling procedure preserve its success
probability and the displayed time and space bounds.
\end{proof}
This is a direct sharpening of a stated consequence, already implicit
in the source's stronger main theorem. It is not the new density
argument of Sections~\ref{sec:checks}--\ref{sec:algorithm}.

\section{Why stronger conclusions do not follow}\label{sec:limits}
\subsection{\texorpdfstring{Sparsification still leaves a \(\log\log n\) factor}{Sparsification still leaves a log log n factor}}
The source sparsification parameter \(\sigma\) yields about
\(2^{O(\sigma n)}\) leaves with occurrence bound
\[
 B=O_k\!\left((\sigma^{-1}\log(2/\sigma))^{k-1}\right).
\]
Even after our improved checks remove the extraneous \(\log m\)
dependence, the resulting exponent tradeoff has the form
\[
 n\left(\sigma+
 \frac{C_k\log(1/\sigma)}{\sqrt{\log n}}\right)
\]
up to additional lower-order logarithms in this estimate.
Its minimizing scale is
\(\sigma=\Theta_k(1/\sqrt{\log n})\), and the result is still
\(\Theta_k(n\log\log n/\sqrt{\log n})\).
This is an obstruction to improving this particular bound by parameter
retuning; it is not a lower bound on all sparsification algorithms.

\subsection{Degree before index compression is not degree after it}
The source's lifted incidence polynomial has degree \(O(n)\).
But its \(O(n\log\log n)\) wires encode subset ranks in binary.
A tensor's degree in those index bits need not equal its degree
in the original incidence indicators.
Using the old \(O(n)\) degree in the compressed network's Hermite
estimate without a new proof would be invalid.

There is a simple dimension obstruction to a universal fix.
If an injective encoding of a finite set \(\Omega\) into
\(\{0,1\}^M\) permits every function on \(\Omega\) to be represented
by a polynomial of degree at most \(K\), then
\[
 |\Omega|\le\sum_{j=0}^{K}\binom Mj.
\]
Multilinearization preserves values and degree on Boolean inputs,
and the right side is the dimension of the representing space.
For all \(d\)-subsets of \(N=\Theta(n(\log n)^c)\) elements, with
fixed \(c>0\) and \(d=\Theta(n)\), the left side has logarithm
\(\Theta(n\log\log n)\).
With \(M=O(n\log\log n)\) and \(K=O(n)\), the right side has
logarithm only \(O(n\log\log\log n)\).
A successful low-degree representation must exploit the special
coefficient functions, not merely relabel arbitrary subset tables.

\subsection{Sharper partition counting alone is insufficient}
The existing mode loop has \(\Theta(M/D_i)\) available modes at each
weight in filled packet families and precision
\(\Theta(MD_i/u)\).
Choose a sufficiently small constant \(c>0\) and weights in
\([J,2J]\), where \(J=\floor{cD_i/\sqrt u}\).
There are \(\Theta(M/\sqrt u)\) such modes.
All their zero--one subsets fit the weighted precision budget after
choosing \(c\) sufficiently small. Their total degree is at most \(M\);
the same construction fits the per-packet row-degree cap.
Thus this particular enumeration can still contain
\(2^{\Omega(M/\sqrt u)}\) tuples after these natural caps.
One must change the way this sum is evaluated to obtain a better order
from this route.

\subsection{A local Cauchy rank obstruction}\label{subsec:smith}
Over \(\Z_{(2)}\), or the source's unramified extension,
consider the odd-unit Cauchy block
\[
 K_{ij}=\frac1{a+2i-2j},\qquad 0\le i,j<N,
\]
where \(a\) is a unit. Its Smith invariant valuations are exactly
\begin{equation}\label{eq:smith}
 e_j=2j+2v_2(j!)=4j-2s_2(j),
\end{equation}
where \(s_2(j)\) is the binary digit sum.

To prove this, the Cauchy determinant identity gives, for an \(r\)-minor
on row and column indices \(x_1,\ldots,x_r\) and \(y_1,\ldots,y_r\),
\[
 v_2(\det K[x,y])
 =r(r-1)+v_2\!\prod_{i<j}(x_j-x_i)
          +v_2\!\prod_{i<j}(y_j-y_i).
\]
The denominators are units.
For arbitrary distinct integer indices, a Vandermonde determinant
is divisible by \(\prod_{j=0}^{r-1}j!\):
its quotient is the integer determinant of the binomial-polynomial
evaluation matrix. Consecutive indices attain equality.
The \(r\)th determinantal ideal therefore has valuation
\[
 r(r-1)+2\sum_{j=0}^{r-1}v_2(j!).
\]
Subtract consecutive values and use Legendre's formula to obtain
\eqref{eq:smith}. The same argument applies in the unramified extension,
where every nonzero element is a power of two times a unit.

Hence the local valuation profile is linear in the rank index.
An integral change to a Newton or factorial basis may improve
constants, but cannot produce a quadratic valuation profile for
these blocks. This says nothing about all SAT algorithms; it rules
out one specific proposed local explanation for a stronger exponent.

\section{Further research questions}\label{sec:research}
The density refinement suggests several concrete next steps.
\begin{enumerate}
\item \textbf{More economical high-rate local codes.}
Replace exhaustive search by an explicit binary construction with
redundancy \(O(d/B)\), relative distance
\(\Omega(1/(B\log(2B)))\), and a substantially smaller graph degree.
This would improve finite thresholds and preprocessing, though it
does not by itself change the theorem's leading logarithmic scale.

\item \textbf{Use the structure of violation vectors.}
Our initial checks reject every small nonzero vector in \(\F_2^m\).
Only vectors realizable as \(v(x)\) for some Boolean assignment matter.
A deterministic construction exploiting this restricted set might
lower row support beyond the code-distance tradeoff. Its construction
cost must include finding the checks, without an oracle for \(\#F\).

\item \textbf{Aggregate shared supports before communicating them.}
The incidence expansion has the form
\[
 \sum_{S} c(S)\,2^{n-|\bigcup S|}.
\]
Its final variable-side weight depends only on the union of the
original variable supports. A local tensor representation communicating
\(O(n)\) bits about shared supports, with degree at most \(n^{9/10}\)
and entry cost \(2^{O(n^{9/10})}\), would remove the remaining
\(\log\log n\) factor in the unrestricted algorithm.

This is an unproved target. Defining
\(C(U)=\sum_{S:\,\cup S=U}c(S)\) is not a solution:
its direct formula is a global subset sum, and broadcasting \(U\)
to every test group restores the repeated incidence cost.

\item \textbf{A recursive degree-sensitive contraction invariant.}
At a single call, if the actual lifted row polynomial has degree \(K\),
the Hermite degree drops from \(M(2L-1)\) to \(K(2L-1)\).
With controlled inherited precision, the mode bound becomes
\(\exp(O(\sqrt{MK/u}))\).
To obtain a new recursive algorithm, this degree bound must survive
both binary index compression and the creation of child tensors.
Neither property follows from the source's incidence-degree estimate.

\item \textbf{Avoid explicit mode enumeration.}
The partition and Smith-profile obstructions above leave open
algebraic aggregation, cancellation, or a new decomposition of the
entire mode sum. Any claimed gain must account for the degree and
entry costs of the new contracted objects.
\end{enumerate}

\section{Verification and deliverables}\label{sec:verification}
The archive contains this note as \texttt{report.tex} and
\texttt{report.pdf}, component-check programs, their output,
a source manifest, and build instructions.

The mathematical proof above is the basis for the asymptotic claims.
Finite tests check implementations of its component formulas and
look for small counterexamples; they do not prove asymptotic complexity.
In particular:
\begin{itemize}
\item The density checks exercise the expander port construction,
entropy and integer parameter inequalities, local-code/Tanner
constraints, and greedy rejection on small examples with a
brute-force counting oracle.
\item The Hermite-retuning checks compare the exact grid
multiplicity, Hermite valuations, and every retained coefficient's
precision inequality on 364 parameter cases, 33,566 retained nodes,
and 434,996 coefficient inequalities.
\item The Smith-profile checker examines all 5,072 minors of
24 rational Cauchy matrices of sizes up to \(6\times6\), using exact
rational arithmetic.
\end{itemize}
The archived logs state the exact checks performed.
No new Lean or Rocq formalization, full fast \(\#k\)-SAT implementation,
or practical benchmark is claimed. The source snapshot is identified
by commit and blob hashes, and a helper can retrieve that exact version
in an environment with network access.

\paragraph{Main outcome.}
The proposed change gives a smaller proved exponent scale for
unbounded, slowly growing clause densities while preserving exactness,
determinism, and polynomial space. Extending the improvement to the
unrestricted worst case remains a separate research problem.

\begingroup\small
\begin{thebibliography}{9}
\bibitem{OAI26}
OpenAI.
\emph{The Exponential Time Hypothesis is False}.
Preprint, 17 October 2026, revised 25 October 2026.
Repository snapshot:
\href{https://github.com/openai/math/tree/42907c9e999cc00d6abe7a50bfc1023250edd1b9/preprints/The-Exponential-Time-Hypothesis-is-False-October-25-2026}{openai/math, commit 42907c9e999c}.
In particular, Theorem 1.1; the sparse initial-check and greedy-hash
lemmas in the main source; and the files
\texttt{quantitative-reduction.tex} and
\texttt{quantitative-contraction.tex} in its \texttt{build} directory.

\bibitem{GG81}
Ofer Gabber and Zvi Galil.
Explicit constructions of linear-sized superconcentrators.
\emph{Journal of Computer and System Sciences} 22(3), 407--420, 1981.
\href{https://doi.org/10.1016/0022-0000(81)90040-4}
{doi:10.1016/0022-0000(81)90040-4}.

\bibitem{Karayel23}
Emin Karayel.
\emph{Expander Graphs}.
Archive of Formal Proofs, 3 March 2023.
\href{https://isa-afp.org/entries/Expander_Graphs.html}
{Project entry} and
\href{https://isa-afp.org/browser_info/current/AFP/Expander_Graphs/Expander_Graphs_MGG.html}
{Expander\_Graphs\_MGG theory};
the graph and bound are \texttt{mgg\_graph} and
\texttt{mgg\_numerical\_radius}.
\end{thebibliography}
\endgroup
\end{document}

